\chapter{Sprint 6: n8n Automation and Alerting}
\chaptertoc

\section{Introduction}
Sprint 6 connected queue metrics to n8n so alerts could leave the backend in a controlled way. The focus was not just sending payloads, but keeping delivery reliable, auditable, and safe to repeat.

\section{Sprint Backlog}
Sprint 6 backlog focused on operational automation:
\begin{itemize}
    \item define and validate a stable webhook payload contract,
    \item implement robust webhook transmission with retries and backup logging,
    \item integrate alert triggers from runtime threshold detection,
    \item implement n8n routing for authorization, cooldown, formatting, and delivery,
    \item provide test scripts for repeated workflow validation.
\end{itemize}

\section{Design}

\subsection{Component Diagram}
\begin{figure}[H]
    \centering
    \fbox{\parbox[c][6cm][c]{0.85\textwidth}{\centering Placeholder for Sprint 6 automation component diagram}}
    \caption{Components involved in backend-to-n8n-to-Telegram automation flow.}
    \label{fig:sprint6_component}
\end{figure}

The automation path has five pieces: queue runtime metrics, webhook payload serializer, HTTP delivery client, n8n workflow router, and notification endpoints. Alert archival is routed back to the backend through a dedicated API endpoint.

\subsection{Data Flow Diagram}
\begin{figure}[H]
    \centering
    \fbox{\parbox[c][6cm][c]{0.85\textwidth}{\centering Placeholder for Sprint 6 automation data flow diagram}}
    \caption{Alert payload flow from backend metrics to external notifications.}
    \label{fig:sprint6_data_flow}
\end{figure}

Data flow starts with periodic metrics dispatch from the runtime loop, branches on alert presence and cooldown in n8n, and then either sends Telegram notifications or returns acknowledgment without external dispatch. Archived alert payloads are persisted through \texttt{/api/alerts/archive}.

\section{Implementation}
\subsection{Webhook Payload Contract}
The webhook schema module defines \texttt{QueuePayload}, a structured contract carrying:
\begin{itemize}
    \item source metadata (timestamp, frame, source, feed ID),
    \item queue metrics ($\lambda$, $\mu$, wait time, stability),
    \item alert flags and severity fields.
\end{itemize}
The module includes payload validation before dispatch.

\subsection{Reliable Webhook Delivery}
The webhook delivery client sends payloads to n8n with timeout, retry policy, optional shared-secret header, and success/failure statistics. Payload persistence is handled through the SQLite database layer, keeping one source of truth for historical records.

\subsection{Runtime Alert Dispatch Integration}
In the main runtime loop, webhook sending is offloaded to \texttt{AsyncWebhookDispatcher} to avoid blocking the perception pipeline. The implementation supports periodic metric dispatch and immediate warning dispatch with deduplication windows to limit repeated alert spam.

\subsection{n8n Workflow Routing and Telegram Delivery}
The n8n alert workflow implements:
\begin{itemize}
    \item webhook reception and secret validation,
    \item alert detection and cooldown gating,
    \item alert history archival via backend API,
    \item message formatting and Telegram send node,
    \item explicit acknowledge/reject branches.
\end{itemize}
That flow keeps the no-alert, alert-suppressed, and alert-delivered branches distinct enough to audit later.

\subsection{Automation Testing Artifacts}
Workflow validation is supported by dedicated PowerShell automation scripts for single-alert checks, automated execution polling, and multi-feed branch scenarios. These scripts automate webhook submission and branch-validation workflows during sprint testing.

\section{Conclusion}
Sprint 6 turned analytics output into actionable notifications and kept the historical trail inside the backend.
